% ==============================================================================
% CONFIG: FORMATTING.TEX
% Cấu hình định dạng trang, tiêu đề chương mục, header/footer, khoảng cách,
% bảng biểu và mã nguồn sử dụng hệ font thống nhất Latin Modern.
% ==============================================================================

% ------------------------------------------------------------------------------
% 1. GIÃN DÒNG VÀ KHOẢNG CÁCH ĐOẠN VĂN
% Line spacing: 1.35 (phù hợp tỷ lệ vàng văn bản học thuật)
% Thụt đầu dòng: 1.0 cm
% Khoảng cách giữa các đoạn văn: 6pt
% ------------------------------------------------------------------------------
\setstretch{1.35}
\setlength{\parindent}{1.0cm}
\setlength{\parskip}{6pt}

% Giới hạn độ sâu phân cấp đánh số và hiển thị mục lục
\setcounter{secnumdepth}{3}
\setcounter{tocdepth}{3}

% ------------------------------------------------------------------------------
% 2. VIỆT HÓA CÁC TÊN TIÊU ĐỀ MẶC ĐỊNH
% ------------------------------------------------------------------------------
\renewcommand{\chaptername}{CHƯƠNG}
\renewcommand{\contentsname}{MỤC LỤC}
\renewcommand{\listfigurename}{DANH MỤC HÌNH ẢNH}
\renewcommand{\listtablename}{DANH MỤC BẢNG}
\renewcommand{\bibname}{TÀI LIỆU THAM KHẢO}
\renewcommand{\figurename}{Hình}
\renewcommand{\tablename}{Bảng}
\renewcommand{\appendixname}{PHỤ LỤC}

% ------------------------------------------------------------------------------
% 3. ĐỊNH DẠNG CLEVEREF CHO TIẾNG VIỆT
% ------------------------------------------------------------------------------
\crefname{figure}{Hình}{Hình}
\Crefname{figure}{Hình}{Hình}
\crefname{table}{Bảng}{Bảng}
\Crefname{table}{Bảng}{Bảng}
\crefname{equation}{Phương trình}{Phương trình}
\Crefname{equation}{Phương trình}{Phương trình}
\crefname{chapter}{Chương}{Chương}
\Crefname{chapter}{Chương}{Chương}
\crefname{section}{Mục}{Mục}
\Crefname{section}{Mục}{Mục}
\crefname{subsection}{Mục}{Mục}
\Crefname{subsection}{Mục}{Mục}
\crefname{subsubsection}{Mục}{Mục}
\Crefname{subsubsection}{Mục}{Mục}
\crefname{appendix}{Phụ lục}{Phụ lục}
\Crefname{appendix}{Phụ lục}{Phụ lục}
\crefname{listing}{Đoạn mã}{Đoạn mã}
\Crefname{listing}{Đoạn mã}{Đoạn mã}

% ------------------------------------------------------------------------------
% 4. ĐỊNH DẠNG TIÊU ĐỀ CHƯƠNG VÀ MỤC (TITLESEC)
% Sử dụng cỡ chữ LaTeX tiêu chuẩn (\large, \Large, \LARGE) và khoảng cách hợp lý.
% ------------------------------------------------------------------------------
% Tiêu đề chương có đánh số
\titleformat{\chapter}[display]
  {\normalfont\bfseries\centering}
  {\large\chaptername\ \thechapter}
  {0.6ex}
  {\Large}

% Tiêu đề chương không đánh số (Mục lục, Lời cảm ơn, v.v.)
\titleformat{name=\chapter,numberless}[display]
  {\normalfont\bfseries\centering}
  {}
  {0ex}
  {\Large}

\titlespacing*{\chapter}{0pt}{10pt}{20pt}

% Mục cấp 1 (Section): 1.1. Tiêu đề
\titleformat{\section}
  {\normalfont\large\bfseries}
  {\thesection.}{0.5em}{}
\titlespacing*{\section}{0pt}{14pt}{6pt}

% Mục cấp 2 (Subsection): 1.1.1. Tiêu đề
\titleformat{\subsection}
  {\normalfont\normalsize\bfseries}
  {\thesubsection.}{0.5em}{}
\titlespacing*{\subsection}{0pt}{10pt}{4pt}

% Mục cấp 3 (Subsubsection): 1.1.1.1. Tiêu đề
\titleformat{\subsubsection}
  {\normalfont\normalsize\itshape}
  {\thesubsubsection.}{0.5em}{}
\titlespacing*{\subsubsection}{0pt}{8pt}{3pt}

% ------------------------------------------------------------------------------
% 5. ĐỊNH DẠNG CAPTION HÌNH ẢNH VÀ BẢNG BIỂU
% Bảng: Caption đặt ở TRÊN bảng
% Hình: Caption đặt ở DƯỚI hình
% ------------------------------------------------------------------------------
\captionsetup{
    font=small,
    labelfont=bf,
    labelsep=colon,
    justification=centering,
    skip=6pt
}

% ------------------------------------------------------------------------------
% 6. MÀU SẮC VÀ ĐỊNH DẠNG HYPERLINK
% Tone màu chuẩn học thuật, nghiêm túc.
% ------------------------------------------------------------------------------
\definecolor{academicblue}{RGB}{10, 45, 95}
\definecolor{codegreen}{RGB}{40, 140, 60}
\definecolor{codegray}{RGB}{110, 110, 110}
\definecolor{codepurple}{RGB}{125, 30, 140}
\definecolor{codebg}{RGB}{248, 249, 250}
\definecolor{framegray}{RGB}{220, 220, 220}

\hypersetup{
    colorlinks=true,
    linkcolor=academicblue,
    citecolor=academicblue,
    urlcolor=academicblue,
    filecolor=academicblue
}

% ------------------------------------------------------------------------------
% 7. HEADER VÀ FOOTER (FANCYHDR)
% ------------------------------------------------------------------------------
\pagestyle{fancy}
\fancyhf{}

% Header: Hiển thị tên chương hiện tại bên trái
\fancyhead[L]{\small\itshape\nouppercase{\leftmark}}
\fancyhead[R]{}
\renewcommand{\headrulewidth}{0.4pt}

% Footer: Số trang ở chính giữa
\fancyfoot[C]{\thepage}
\renewcommand{\footrulewidth}{0pt}

\setlength{\headheight}{16pt}

% Trang đầu tiên của chương dùng kiểu 'plain' (chỉ có số trang ở footer)
\fancypagestyle{plain}{
    \fancyhf{}
    \fancyfoot[C]{\thepage}
    \renewcommand{\headrulewidth}{0pt}
    \renewcommand{\footrulewidth}{0pt}
}

% ------------------------------------------------------------------------------
% 8. ĐỊNH DẠNG MÃ NGUỒN (LISTINGS) VỚI LATIN MODERN MONO
% ------------------------------------------------------------------------------
\lstdefinestyle{pythonstyle}{
    language=Python,
    backgroundcolor=\color{codebg},
    commentstyle=\color{codegreen}\itshape,
    keywordstyle=\color{academicblue},
    numberstyle=\tiny\color{codegray},
    stringstyle=\color{codepurple},
    basicstyle=\ttfamily\footnotesize,
    breakatwhitespace=false,
    breaklines=true,
    captionpos=b,
    keepspaces=true,
    numbers=left,
    numbersep=8pt,
    showspaces=false,
    showstringspaces=false,
    showtabs=false,
    tabsize=4,
    frame=single,
    rulecolor=\color{framegray},
    xleftmargin=12pt,
    framexleftmargin=12pt
}
\lstset{style=pythonstyle}
